\documentclass[11pt,a4paper]{article}
\usepackage[T1]{fontenc}
\usepackage[utf8]{inputenc}
\usepackage{amsmath,amsthm,mathtools}
\usepackage{newtxtext,newtxmath}
\usepackage[margin=27mm,headheight=15pt]{geometry}
\usepackage{microtype}
\usepackage{xcolor,booktabs,array,enumitem}
\usepackage{fancyhdr}
\usepackage[most]{tcolorbox}
\usepackage{listings}
\usepackage{hyperref}
\usepackage[nameinlink,noabbrev]{cleveref}
\definecolor{Ink}{HTML}{20344B}
\definecolor{Accent}{HTML}{176B78}
\definecolor{Pale}{HTML}{F0F5F7}
\hypersetup{colorlinks=true,linkcolor=Accent,citecolor=Accent,urlcolor=Accent,
 pdftitle={Every ordinal as a point-separating game value},
 pdfauthor={Research manuscript prepared for Vladimir Reshetnikov}}
\pagestyle{fancy}
\fancyhf{}
\fancyhead[L]{\small\textsc{Lexicographic cubes and ordinal game lengths}}
\fancyhead[R]{\small\thepage}
\renewcommand{\headrulewidth}{0.3pt}
\setlength{\parskip}{4pt}
\setlength{\parindent}{1.25em}
\setlist{itemsep=2pt,topsep=4pt}
\newtheorem{theorem}{Theorem}[section]
\newtheorem{lemma}[theorem]{Lemma}
\newtheorem{proposition}[theorem]{Proposition}
\newtheorem{corollary}[theorem]{Corollary}
\theoremstyle{definition}
\newtheorem{definition}[theorem]{Definition}
\newtheorem{example}[theorem]{Example}
\theoremstyle{remark}
\newtheorem{remark}[theorem]{Remark}
\newcommand{\ps}{\operatorname{ps}}
\newcommand{\sep}{\operatorname{sep}}
\newcommand{\lex}{\mathrm{lex}}
\newcommand{\fin}{\mathrm{fin}}
\newcommand{\dd}{\downarrow}
\newcommand{\restr}{\mathbin{\upharpoonright}}
\newcommand{\cat}{\mathbin{\frown}}
\newcommand{\two}{\{0,1\}}
\newcommand{\cfr}{\mathfrak c}
\newcommand{\Pow}{\mathcal P}
\newcommand{\B}{\mathbb B}
\newcommand{\N}{\mathbb N}
\newcommand{\tauD}{\tau_{\dd}}
\newcommand{\tauF}{\tau_{\fin}}
\newtcolorbox{resultbox}{colback=Pale,colframe=Accent,boxrule=0.6pt,
 arc=1mm,left=3mm,right=3mm,top=2mm,bottom=2mm,breakable}
\lstset{basicstyle=\small\ttfamily,columns=fullflexible,keepspaces=true,
 breaklines=true,frame=single,rulecolor=\color{Accent},showstringspaces=false}
\title{\color{Ink}\LARGE Every ordinal as a point-separating game value\\[5pt]
\large Lexicographic cubes, compact $T_1$ realizations,\\
and a proposed solution of a recent spectrum problem}
\author{Research manuscript prepared for Vladimir Reshetnikov}
\date{20 September 2026}
\begin{document}
\maketitle
\begin{abstract}
For every ordinal $\lambda$, let $\B_\lambda=\two^\lambda$ carry its
lexicographic order. We prove that its topology of all downward-closed sets has
point-separating game value exactly $\lambda$. The lower bound has a particularly
short strategy: after each query, retain an entire immediate successor cylinder.
At limit stages, take the union of the retained prefixes. We then prove the
same exact value for the topology generated by downward-closed sets with finitely
many points deleted. This second topology is compact and $T_1$. Its lower bound
uses a finite-prefix homogeneity lemma and transfinite fusion in blocks of length
$\omega$. Consequently every ordinal occurs as the point-separating value of a
compact $T_1$ space. In particular, the value $\omega^2+1$ gives an affirmative
answer to Problem~1.2 of Chiozini, Csern\'ak, and Soukup, in its stated generality.
The infinite examples are not Hausdorff. We also give an order-embedding
characterization of the invariant for downward topologies, exact nonadaptive
separation numbers in the countable case, and explicit winning strategies for
both players on the appropriate sides of the threshold.
\end{abstract}
\begin{resultbox}
\textbf{Status and scope.} This is a proposed research solution with full proofs,
not an independently refereed or machine-formalized result. The checked source is
arXiv:2510.05754v3. The construction answers its unrestricted topological-space
question and even yields compact $T_1$ witnesses. It does \emph{not} establish
realization by Hausdorff, regular Hausdorff, or zero-dimensional Hausdorff spaces.
Throughout this manuscript, compactness means the finite-subcover property;
Hausdorffness is a separate condition.
\end{resultbox}
\clearpage
\tableofcontents
\clearpage

\section{The selected problem and the main result}

\subsection{A concrete ordinal-spectrum question}
Chiozini, Csern\'ak, and Soukup introduce an ordinal-valued point-separating game
invariant in \cite{CCS}. Their Problem~1.2 asks whether a topological space can
satisfy
\begin{equation}\label{eq:problem}
                 \omega^2<\ps(X)<\omega_1.
\end{equation}
The question is posed in a framework allowing $T_0$ spaces. The same source
constructs zero-dimensional Hausdorff examples for all values at most
$\omega^2$. The version checked here, v3, was submitted on 29 May 2026; its PDF
is dated 1 June 2026 and still explicitly poses the question on pages 4 and 16.
The bibliographic record retains the earlier title \emph{Cut-and-choose games
in topological spaces}, while the v3 PDF has the title recorded in our reference.

Our approach replaces complicated topological constructions by an ordered
search space. A binary sequence indexed by an ordinal can be read one coordinate
at a time. The point is to choose a topology in which no alternative query can
uniformly finish earlier. Downward-closed queries have exactly the required
property. Finite deletions permit a $T_1$ refinement without destroying the
transfinite obstruction.

\subsection{The construction}
For an ordinal $\lambda$, define
\[
 \B_\lambda=\{x:x\colon\lambda\longrightarrow\two\}.
\]
If $x\ne y$, let $\xi$ be their first coordinate of disagreement. Set
$x<_{\lex}y$ when $x(\xi)=0$ and $y(\xi)=1$. This is a linear order because
$\lambda$ is well ordered.

On this ordered set, use the following two topologies:
\begin{align}
 \tauD&=\{D\subseteq\B_\lambda:
       y\leq_{\lex}x\in D\Longrightarrow y\in D\},\label{eq:down}\\
 \tauF&=\left\{\bigcup_{i\in I}(D_i\setminus F_i):
       D_i\in\tauD,\ F_i\subseteq\B_\lambda\text{ finite}\right\}.
       \label{eq:finite-top}
\end{align}
Write $X_\lambda^\dd=(\B_\lambda,\tauD)$ and
$X_\lambda^\fin=(\B_\lambda,\tauF)$.

\begin{theorem}[Ordinal realization]\label{thm:main}
For every ordinal $\lambda$,
\begin{equation}\label{eq:main}
             \ps(X_\lambda^\dd)=\ps(X_\lambda^\fin)=\lambda.
\end{equation}
The space $X_\lambda^\dd$ is compact and $T_0$, and
$X_\lambda^\fin$ is compact and $T_1$. For either space, Seeker has a winning
strategy in the length-$\beta$ point-separating game exactly when
$\beta\geq\lambda$; Hider has a winning strategy exactly when $\beta<\lambda$.
For infinite $\lambda$, neither construction is Hausdorff.
\end{theorem}

\begin{corollary}\label[corollary]{cor:answer}
There is a compact $T_1$ space of cardinality $\cfr=2^{\aleph_0}$ satisfying
$\ps(X)=\omega^2+1$. More generally, every countably infinite ordinal occurs
as $\ps(X)$ for a compact $T_1$ space of cardinality $\cfr$.
\end{corollary}
\begin{proof}
Take $X=X_{\omega^2+1}^\fin$, or $X=X_\lambda^\fin$ for an arbitrary countably
infinite $\lambda$. The underlying set has cardinality $2^{|\lambda|}=\cfr$.
Apply \cref{thm:main}.
\end{proof}

No continuum hypothesis, forcing assumption, or large-cardinal hypothesis is
used. The theorem is formulated for all set ordinals, not only countable ones.
Its primary claim of research interest is the realization statement in relation
to \eqref{eq:problem}; the elementary order-theoretic consequences below are not
claimed to be historically new in isolation.

\section{Game conventions and ordinal preliminaries}

\subsection{The point-separating game}
We use the point-separating game of \cite[Definition~1.1]{CCS}.
Fix a topological space $X$ and an ordinal $\beta$. At turn $\xi<\beta$, Seeker
chooses an open set $U_\xi$, and Hider answers $i_\xi\in\two$. Put
\[
 V_\xi=\begin{cases}U_\xi&i_\xi=1,\\X\setminus U_\xi&i_\xi=0,
 \end{cases}
 \qquad
 W_\eta=\bigcap_{\xi<\eta}V_\xi\quad(\eta\leq\beta),
\]
with $W_0=X$. Seeker wins if $|W_\beta|\leq1$; otherwise Hider wins. The
point-separating number is
\[
 \ps(X)=\min\{\beta:\text{Seeker has a winning strategy in }G_\beta(X)\}.
\]
A strategy may depend on the complete earlier history, including at limit
stages. No computability restriction is imposed.

If the candidate set ever becomes empty, it remains empty, and Seeker wins.
Thus a lower bound must preserve actual points through limit intersections.
Showing only that every finite position is nonempty is insufficient. All lower
bounds in this manuscript maintain an explicit full cylinder inside the candidate
set, or a specified finite set at the very end.

There are two distinct logical assertions worth keeping separate:
\[
 \text{Seeker has no winning strategy},\qquad
 \text{Hider has a winning strategy}.
\]
The second implies the first, but no general determinacy principle is being
assumed. We prove the stronger assertion directly for our spaces. The strategies
therefore also refute every fixed Seeker strategy in the interpretation where a
hidden point is chosen in advance: the resulting final candidate set contains
two points, either of which could have generated all the recorded answers.

\subsection{Ordinal addition is not cardinal arithmetic}
All additions and multiplications of ordinals below are the ordinary ordinal
operations, not the natural operations. In particular, $\omega\cdot\alpha$
means $\alpha$ consecutive blocks of order type $\omega$.

Every ordinal $\lambda$ has a unique decomposition
\begin{equation}\label{eq:division}
                    \lambda=\omega\cdot\alpha+n,
           \qquad n<\omega.
\end{equation}
This is division by $\omega$ on the left. One way to obtain it is to take the
largest $\alpha$ with $\omega\cdot\alpha\leq\lambda$. Such an $\alpha$ exists
because multiplication by $\omega$ on the left is increasing and continuous in
the right argument, and $\omega\cdot\alpha\geq\alpha$ bounds the relevant
set of candidates. The remaining interval has order type less than $\omega$.
Uniqueness follows from
$\omega\cdot\alpha+n<\omega\cdot(\alpha+1)$.

At a limit ordinal $\delta$,
\begin{equation}\label{eq:continuity}
             \sup_{\xi<\delta}\omega\cdot\xi=\omega\cdot\delta.
\end{equation}
For a fixed $\xi$, a strictly increasing sequence of finite integers $k_j$
satisfies
\[
 \sup_{j<\omega}(\omega\cdot\xi+k_j)=\omega\cdot(\xi+1).
\]
These two observations account for all limit-stage bookkeeping in the $T_1$
proof. Merely enumerating a countable ordinal by $\N$ would not preserve this
order of dependence.

\section{Lexicographic cylinders and the two topologies}

\subsection{Prefix geometry}
For $\eta\leq\lambda$ and $s\in\two^\eta$, define the full cylinder
\begin{equation}\label{eq:cylinder}
        C_s=\{x\in\B_\lambda:x\restr\eta=s\}.
\end{equation}
The length of $s$ is its ordinal domain, denoted $\operatorname{lh}(s)$.
The empty prefix has cylinder $\B_\lambda$. A prefix of length $\lambda$ has a
singleton cylinder. Every shorter prefix has at least two extensions.

\begin{lemma}[Cylinder geometry]\label[lemma]{lem:cylinders}
The cylinders have the following properties.
\begin{enumerate}[label=(\alph*)]
\item Each $C_s$ is nonempty and convex in the lexicographic order.
\item If $\operatorname{lh}(s)<\lambda$, then
$C_s=C_{s\cat0}\mathbin{\dot\cup}C_{s\cat1}$, and every element of the first
child is less than every element of the second.
\item For an increasing compatible family of prefixes $(s_\xi)_{\xi<\delta}$,
with union $s$, one has
$\bigcap_{\xi<\delta}C_{s_\xi}=C_s$.
\end{enumerate}
\end{lemma}
\begin{proof}
A prefix extends to a point of $\B_\lambda$ by assigning zero to all remaining
coordinates. If two extensions of the same prefix bracket a third point, that
point cannot first disagree with the prefix before its end, proving convexity.
The second statement follows by inspection of the next bit. For the last
statement, belonging to every cylinder is exactly the condition of agreeing
with every $s_\xi$, hence with their union.
\end{proof}

Notice that these are \emph{full} cylinders: no restrictions are placed on later
coordinates. This is what prevents a hidden loss of choices at a limit.

\subsection{The downward topology}
The family $\tauD$ is closed under arbitrary unions and arbitrary intersections,
and contains the empty set and the whole space. It is therefore a topology,
indeed an Alexandrov topology. This use of ``Alexandrov'' means closure under
arbitrary intersections of open sets; it does not refer to the double-arrow
space.

For $x<_{\lex}y$, the principal initial segment
$\{z:z\leq_{\lex}x\}$ contains $x$ but not $y$, so $X_\lambda^\dd$ is $T_0$.
Its greatest point is $1^\lambda$. The only downward-closed open set containing
this point is the whole space, so every open cover contains the whole space as
one of its members. In particular, $X_\lambda^\dd$ is compact.

If $\lambda>0$, every nonempty downward-open set contains $0^\lambda$. Thus
$X_\lambda^\dd$ is not $T_1$ and is not Hausdorff. This simple topology will first
provide a transparent proof of the ordinal obstruction.

\subsection{The finite-deletion refinement}
The sets $D\setminus F$ in \eqref{eq:finite-top} form a base because
\[
 (D\setminus F)\cap(E\setminus G)
             =(D\cap E)\setminus(F\cup G).
\]
This is the smallest topology containing the downward topology and all cofinite
sets.

\begin{proposition}\label[proposition]{prop:compact}
For every ordinal $\lambda$, $X_\lambda^\fin$ is compact and $T_1$. It is discrete
when $\lambda$ is finite and is not Hausdorff when $\lambda$ is infinite.
\end{proposition}
\begin{proof}
For every $x$, the set $\B_\lambda\setminus\{x\}$ is basic open, so all singletons
are closed. This proves $T_1$.

Let an open cover be given, and choose a member $U$ containing $1^\lambda$.
There is a basic set $D\setminus F$ with
$1^\lambda\in D\setminus F\subseteq U$. Downward closure forces
$D=\B_\lambda$, so $U$ misses only finitely many points. At most one further
cover member per missing point gives a finite subcover.

When $\lambda$ is finite, the underlying set is finite and every subset is
cofinite, so the topology is discrete. Now suppose $\lambda\geq\omega$. Let
$y$ have first bit $1$ and all later bits $0$, and let $t=1^\lambda$. Then
$y\ne t$, and infinitely many points lie below $y$. Every neighborhood of $y$
contains all those predecessors with at most finitely many exceptions. Every
neighborhood of $t$ is cofinite, by the preceding argument. These two types of
neighborhoods must intersect, so $y$ and $t$ have no disjoint neighborhoods.
\end{proof}

Neither topology should be confused with the product topology on $\two^\lambda$
or with its two-sided lexicographic order topology. The exact collection of
legal queries is essential to the result.

\section{An exact value for downward queries}

\subsection{Seeker's coordinate strategy}
For a prefix $s$ of length $\eta<\lambda$, define
\begin{equation}\label{eq:seeker-cut}
 D_s=\{x:x\restr\eta<_{\lex}s\}
      \ \cup\ \{x:x\restr\eta=s\text{ and }x(\eta)=0\}.
\end{equation}
This is a downward-closed subset of the \emph{whole} ordered space, and
\begin{equation}\label{eq:cut-trace}
                      D_s\cap C_s=C_{s\cat0}.
\end{equation}

\begin{proposition}[Upper bound]\label[proposition]{prop:upper}
Seeker wins $G_\lambda(X_\lambda^\dd)$ and $G_\lambda(X_\lambda^\fin)$.
\end{proposition}
\begin{proof}
Start with the empty prefix. At turn $\eta$, query $D_s$, where $s$ is the
prefix determined by the previous answers. An answer $1$ determines the next
bit to be $0$, and an answer $0$ determines it to be $1$. Thus the candidate set
is exactly the new cylinder. At a limit, take the union of the earlier prefixes;
\cref{lem:cylinders}(c) gives the same invariant. After $\lambda$ turns the
prefix has length $\lambda$, so exactly one point remains. Each query belongs
to both topologies.
\end{proof}

In particular, the strategy does not ask the globally defined question
``is coordinate $\eta$ zero?'' without qualification. That set need not be
open. Instead it uses the global initial segment \eqref{eq:seeker-cut}, whose
trace on the currently known cylinder is precisely the desired question.

\subsection{Hider's one-bit strategy}
\begin{lemma}[An untouched child]\label[lemma]{lem:child}
Let $D$ be downward closed and let $s$ be a prefix of length less than $\lambda$.
Then either
\[
                 C_{s\cat0}\subseteq D
   \quad\text{or}\quad
                 C_{s\cat1}\subseteq\B_\lambda\setminus D.
\]
\end{lemma}
\begin{proof}
If the first inclusion fails, choose $a\in C_{s\cat0}\setminus D$. Every
$b\in C_{s\cat1}$ is greater than $a$. If such a $b$ belonged to $D$, downward
closure would put $a$ in $D$, a contradiction.
\end{proof}

\begin{theorem}[Downward-query threshold]\label{thm:down-exact}
For $\beta<\lambda$, Hider has a winning strategy in
$G_\beta(X_\lambda^\dd)$. Consequently $\ps(X_\lambda^\dd)=\lambda$.
\end{theorem}
\begin{proof}
Maintain a prefix $s_\eta$ of length exactly $\eta$ with
\begin{equation}\label{eq:down-invariant}
                            C_{s_\eta}\subseteq W_\eta.
\end{equation}
Initially this holds for the empty prefix. When Seeker asks $D$ at stage $\eta$,
use the following rule:
\[
 \begin{array}{c|c|c}
 \text{condition}&\text{answer}&\text{retained prefix}\\\hline
 C_{s_\eta\cat0}\subseteq D&1&s_\eta\cat0\\
 C_{s_\eta\cat0}\nsubseteq D&0&s_\eta\cat1
 \end{array}
\]
The second case is legitimate by \cref{lem:child}. In either case the new
cylinder is contained in the old candidate set and in the chosen side of the
query.

At a limit $\eta$, define $s_\eta=\bigcup_{\xi<\eta}s_\xi$.
Its length is $\eta$, and \cref{lem:cylinders}(c) proves
\eqref{eq:down-invariant}; in particular the intersection is not empty.
At the end, $C_{s_\beta}$ still has a free coordinate because $\beta<\lambda$.
It contains two distinct points, so Hider wins. The upper bound is
\cref{prop:upper}.
\end{proof}

This proof is entirely online: Hider uses only the current query and the
retained prefix. It does not inspect a future decision tree or invoke a
determinacy theorem. It already proves \cref{cor:answer} with ``compact $T_0$''
in place of ``compact $T_1$''.

\section{Preserving the value under the \texorpdfstring{$T_1$}{T1} refinement}

A finite deletion can destroy the untouched-child argument. For example, a
query might delete one point from each immediate child. Thus the preceding
proof cannot simply be reused after refining the topology. The replacement
allows finitely many prefix bits to be spent on a single move, while ensuring
that only one block of $\omega$ bits is spent over $\omega$ moves.

\subsection{Finite-prefix homogeneity}
\begin{lemma}[Homogeneous finite extension]\label[lemma]{lem:homogeneous}
Let $s$ have length $\theta$, where $\theta+\omega\leq\lambda$, and let
$U\in\tauF$. There is a nonempty finite binary word $v$ such that
\begin{equation}\label{eq:homogeneous}
       C_{s\cat v}\subseteq U
       \quad\text{or}\quad
       C_{s\cat v}\subseteq\B_\lambda\setminus U.
\end{equation}
\end{lemma}
\begin{proof}
First suppose every point of $U\cap C_s$ has zero in all of the next $\omega$
coordinates. Then $C_{s\cat1}$ is disjoint from $U$, and we are done.

Otherwise choose $x\in U\cap C_s$ with a $1$ among the next $\omega$ bits. Let
$j<\omega$ be the first such position. Thus the next $j$ bits are zero and
$x(\theta+j)=1$. Choose a basic neighborhood
\[
                       x\in D\setminus F\subseteq U,
\]
where $D$ is downward closed and $F$ is finite. Put
$p=s\cat0^{j+1}$. Every point of $C_p$ is strictly below $x$, so
$C_p\subseteq D$.

Choose a finite $k$ with $2^k>|F|$. The $2^k$ cylinders
\[
                         C_{p\cat w}\qquad(w\in\two^k)
\]
are pairwise disjoint and nonempty. At most $|F|$ of them meet $F$. Choose one
that does not, say $C_{p\cat w}$. It lies in $D\setminus F\subseteq U$.
The word $v=0^{j+1}\cat w$ is finite and nonempty. All coordinates used exist
because $\theta+j+1+k<\theta+\omega\leq\lambda$.
\end{proof}

\begin{remark}[Making the rule a strategy]\label[remark]{rem:shortlex}
Among all nonempty finite words satisfying \eqref{eq:homogeneous}, select one
of shortest length and then the lexicographically first of that length. This
is a definite choice rule depending only on $s$ and $U$. Answer according to
the side containing its cylinder. The lemma guarantees that the rule is always
defined when a full $\omega$-block remains. Testing the inclusion of an infinite
cylinder in an arbitrary open set is not claimed to be an effective computation;
strategies here are set-theoretic functions, as in the original game.
\end{remark}

\subsection{Fusion through one block and through limit blocks}
Suppose a full cylinder with prefix length $\omega\cdot\xi$ is contained in the
candidate set at game time $\omega\cdot\xi$, and
$\omega\cdot(\xi+1)\leq\lambda$. Apply \cref{lem:homogeneous} at each of the
next finitely indexed turns. After $r<\omega$ turns of this block, the tracked
prefix has length
\begin{equation}\label{eq:block-length}
                     \omega\cdot\xi+k_r,
     \qquad k_r<\omega,\quad k_0=0,\quad k_{r+1}>k_r.
\end{equation}
The applicability condition remains true: adding $\omega$ to any length in
\eqref{eq:block-length} gives $\omega\cdot(\xi+1)$.

At the end of the block, the union of the prefixes has length exactly
$\omega\cdot(\xi+1)$, because the strictly increasing finite sequence $(k_r)$
is unbounded in $\omega$. Its full cylinder survives. At a limit $\delta$ of
blocks, take the union of all earlier tracked prefixes. Equation
\eqref{eq:continuity} shows that the resulting length is $\omega\cdot\delta$,
and \cref{lem:cylinders}(c) again shows that its full cylinder survives.

The construction therefore satisfies two different invariants. At block
boundaries, game time and prefix length agree. Within a block, prefix length
may run finitely far ahead of game time. It can never reach the end of that
$\omega$-block after only finitely many moves.

\subsection{The finite terminal block}
\begin{lemma}[Finite survival]\label[lemma]{lem:finite}
If a candidate set contains a specified set $A$ of $2^n$ points, Hider can retain
at least $2^{n-r}$ of these points after any $r\leq n$ further binary queries,
regardless of which subsets are allowed as queries.
\end{lemma}
\begin{proof}
For each query retain the larger of its two intersections with the current
tracked finite set. A set of at least $2^m$ points has one side with at least
$2^{m-1}$ points. Induction proves the assertion.
\end{proof}

\begin{theorem}[Finite-deletion threshold]\label{thm:fin-exact}
For $\beta<\lambda$, Hider has a winning strategy in
$G_\beta(X_\lambda^\fin)$. Hence $\ps(X_\lambda^\fin)=\lambda$.
\end{theorem}
\begin{proof}
Write $\lambda=\omega\cdot\alpha+n$ with $n<\omega$. During the first
$\omega\cdot\alpha$ moves, use the homogeneous-extension strategy and the
fusion rule just described. At time $\omega\cdot\xi$ for $\xi\leq\alpha$, it
preserves a full cylinder with prefix length $\omega\cdot\xi$.

Consider any shorter stopping time
$\beta=\omega\cdot\delta+r<\lambda$, where $r<\omega$. Necessarily
$\delta\leq\alpha$. If $\delta<\alpha$, only finitely many moves have been
made in block $\delta$. The tracked prefix has length
$\omega\cdot\delta+k$ for some finite $k$. At least the rest of that infinite
block is free, so its cylinder contains infinitely many points.

If $\delta=\alpha$, then $r<n$. At time $\omega\cdot\alpha$, the surviving
full cylinder has exactly $2^n$ points, indexed by the remaining $n$ bits.
Apply \cref{lem:finite}. After $r$ more moves at least $2^{n-r}\geq2$ points
remain. This also handles $\alpha=0$, when the whole space is finite.

Thus the specified Hider strategy wins at every $\beta<\lambda$.
\Cref{prop:upper} supplies the matching Seeker strategy at $\lambda$.
\end{proof}

\begin{proof}[Completion of the proof of \cref{thm:main}]
The equalities follow from \cref{thm:down-exact,thm:fin-exact}. The separation and
compactness properties were established in \cref{prop:compact} and the preceding
subsection on downward topologies. Seeker's strategy at $\lambda$ extends to any
longer game by making arbitrary additional queries: the candidate set cannot
grow. The given Hider strategies apply at every shorter length. Both players
cannot have winning strategies in the same game, since playing them against each
other produces a single play with only one winner.
\end{proof}

\section{The first example beyond \texorpdfstring{$\omega^2$}{omega squared}, in full detail}

Take
\begin{equation}\label{eq:concrete}
 X=\left(\two^{\omega^2+1},\tauF\right).
\end{equation}
The coordinates before $\omega^2$ are arranged as
\[
 [0,\omega),\ [\omega,\omega\cdot2),\ [\omega\cdot2,\omega\cdot3),\ldots,
\]
followed by the single coordinate $\omega^2$. This arrangement is an ordinal
sequence of blocks, not an arbitrary enumeration of a countable set.

Seeker can read the first block in $\omega$ moves, the second in the next
$\omega$ moves, and so on, using \eqref{eq:seeker-cut}. At time $\omega^2$,
all coordinates below $\omega^2$ have been determined. One further initial-segment
query determines the last bit. This gives $\ps(X)\leq\omega^2+1$.

Against any Seeker play of length $\omega^2$, Hider applies the finite-prefix
homogeneity lemma during the first block and retains a full cylinder with an
$\omega$-long prefix. Repeating the procedure, at time $\omega\cdot m$ a full
cylinder with prefix length $\omega\cdot m$ survives. At time $\omega^2$, the
union prefix $s$ has length $\omega^2$ and
\[
                          \{s\cat0,s\cat1\}\subseteq W_{\omega^2}.
\]
The last bit has not been spent. Therefore Hider wins $G_{\omega^2}(X)$.
The same strategy also wins at every earlier stopping time, so
\[
                             \ps(X)=\omega^2+1.
\]

This witness is compact and $T_1$ by \cref{prop:compact}, has size $\cfr$, and is
not Hausdorff. In the downward topology, an even simpler account is available:
Hider retains a cylinder with prefix length exactly equal to game time, so the
two final extensions survive directly by \cref{thm:down-exact}.

\begin{table}[ht]
\centering
\begin{tabular}{@{}lll@{}}
\toprule
Target value $\lambda$ & Underlying ordered set & $\omega$-block decomposition\\
\midrule
$n<\omega$ & $\two^n$ & $\omega\cdot0+n$\\
$\omega+3$ & $\two^{\omega+3}$ & $\omega\cdot1+3$\\
$\omega^2$ & $\two^{\omega^2}$ & $\omega\cdot\omega$\\
$\omega^2+1$ & $\two^{\omega^2+1}$ & $\omega\cdot\omega+1$\\
$\omega^3+7$ & $\two^{\omega^3+7}$ & $\omega\cdot\omega^2+7$\\
$\omega^\omega$ & $\two^{\omega^\omega}$ & $\omega\cdot\omega^\omega$\\
\bottomrule
\end{tabular}
\caption{For every row, either of the two specified topologies has value
$\lambda$; the finite-deletion topology is compact and $T_1$. The last equality
uses $1+\omega=\omega$ in the exponent of an ordinal power.}
\label{tab:examples}
\end{table}

\section{An order-theoretic characterization}

The same game describes exactly the least length of a lexicographic binary code
for an arbitrary linear order. This gives a direct connection with order types,
rather than only with topological invariants.

For a linear order $L$, let $L^\dd$ denote $L$ with all downward-closed sets open.
Its game value exists: well-order the set of principal initial segments and ask
all of them. These sets separate every pair of distinct points.

\begin{theorem}[Lexicographic coding]\label{thm:embedding}
For every linear order $L$ and every ordinal $\beta$, the following are
equivalent:
\begin{enumerate}[label=(\roman*)]
\item Seeker has a winning strategy in $G_\beta(L^\dd)$.
\item There is an order embedding $L\hookrightarrow(\B_\beta,<_{\lex})$.
\end{enumerate}
Consequently
\begin{equation}\label{eq:embedding-rank}
 \ps(L^\dd)=\min\{\beta:L\text{ order-embeds into }\B_\beta\}.
\end{equation}
\end{theorem}
\begin{proof}
Suppose first that a winning strategy is given. For each $x\in L$, follow the
play in which Hider always chooses the side containing $x$. Encode the answer
``inside'' by $0$ and ``outside'' by $1$, obtaining a word $c(x)\in\two^\beta$.
Distinct points have distinct words, since otherwise both would survive the
same complete play.

Let $x<y$. At their first different answer, the preceding histories agree, so
the strategy asks the same downward-closed set $D$. It is impossible that
$x\notin D$ and $y\in D$. Therefore $x$ answers inside and $y$ answers outside.
In our coding convention the first different bits are $0$ and $1$, respectively.
Hence $c(x)<_{\lex}c(y)$, and $c$ is an order embedding.

Conversely, let $e:L\hookrightarrow\B_\beta$ be an order embedding. For a
current code prefix $s$ of length $\eta<\beta$, ask $e^{-1}(D_s)$, where $D_s$
is the initial segment \eqref{eq:seeker-cut} in $\B_\beta$. This preimage is
downward closed in $L$. The response determines one more code bit. After
$\beta$ moves, all remaining points have the same code, and injectivity of $e$
leaves at most one. Impossible code histories simply give an empty candidate
set, which also wins for Seeker.
\end{proof}

\begin{corollary}[Strict hierarchy of order types]\label[corollary]{cor:strict}
For any ordinals $\alpha,\beta$,
\[
           \B_\alpha\text{ order-embeds into }\B_\beta
                     \quad\Longleftrightarrow\quad\alpha\leq\beta.
\]
\end{corollary}
\begin{proof}
If $\alpha\leq\beta$, extend each $\alpha$-sequence by zeros to length $\beta$;
its first disagreement with another sequence is unchanged. Conversely, an
embedding with $\beta<\alpha$ would give a winning length-$\beta$ strategy on
$X_\alpha^\dd$ by \cref{thm:embedding}, contradicting
\cref{thm:down-exact}.
\end{proof}

Thus the countably infinite lexicographic cubes form a strictly increasing
order-embedding hierarchy, even though all have the same cardinality $\cfr$.
Their game values measure this order-sensitive coding length. A bijection
between the coordinate sets of two countable ordinals does not in general
preserve their lexicographic order.

\section{Adaptive versus nonadaptive separation}

For a $T_0$ space $X$, define
\[
 \sep(X)=\min\bigl\{|\mathcal U|:\mathcal U\subseteq\tau_X,
  \ \forall x\ne y\ \exists U\in\mathcal U\ (|U\cap\{x,y\}|=1)\bigr\}.
\]
Unlike $\ps$, this is a cardinal: all queries in $\mathcal U$ are chosen in
advance. For infinite spaces, this agrees with the $T_0$-pseudoweight discussed
in \cite{CCS}; the definition above avoids finite-cardinal convention issues.
Our calculations here are proved directly.

\subsection{Adjacent pairs force one cut each}
Let
\[
              \two^{<\lambda}=\bigcup_{\eta<\lambda}\two^\eta
\]
be the set of all proper-length prefixes, with sequences of different lengths
regarded as different elements.

\begin{proposition}\label[proposition]{prop:sep-down}
For every ordinal $\lambda$,
\begin{equation}\label{eq:sep-down}
                      \sep(X_\lambda^\dd)=|\two^{<\lambda}|.
\end{equation}
In particular this equals $2^n-1$ for $\lambda=n<\omega$, equals $\aleph_0$
for $\lambda=\omega$, and equals $\cfr$ for
$\omega<\lambda<\omega_1$.
\end{proposition}
\begin{proof}
The family $(D_s)_{s\in\two^{<\lambda}}$ separates all pairs: use their first
coordinate of disagreement. This gives the upper bound.

For each proper prefix $s$ of length $\eta$, form the pair
\[
 a_s=s\cat0\cat1^{[\eta+1,\lambda)},\qquad
 b_s=s\cat1\cat0^{[\eta+1,\lambda)}.
\]
Here the last notation means assigning the indicated constant value to all
remaining coordinates. The point $a_s$ is the maximum of $C_{s\cat0}$, and
$b_s$ is the minimum of $C_{s\cat1}$. They are adjacent: no point lies strictly
between them. Different prefixes give different adjacent pairs.

A downward-closed set can separate at most one adjacent pair in a linear order.
Indeed, if the first pair occurs before the second, containing the lower point
of the second pair forces containment of the upper point of the first pair.
Every separating family must separate all these pairs, proving the lower bound.
The stated special cases follow by counting the levels of the prefix tree.
\end{proof}

\subsection{Finite deletions can split only countably many ordered fibres}
The following lemma handles arbitrary unions of finite-deletion basic sets;
it does not incorrectly assume that an arbitrary open set itself has only
finitely many holes.

\begin{lemma}[Split fibres]\label[lemma]{lem:split-fibres}
Let $L$ be a linear order with the topology generated by downward-closed sets
minus finite sets. Suppose $(A_t)_{t\in I}$ is a family of nonempty subsets
indexed by a linear order, with every point of $A_s$ below every point of $A_t$
whenever $s<t$. For an open set $U$, let
\[
 E=\{t:\varnothing\ne U\cap A_t\ne A_t\}.
\]
Then $E$ is at most countable. More precisely, every element of $E$ has only
finitely many predecessors in $E$.
\end{lemma}
\begin{proof}
Fix $t\in E$, choose $b\in U\cap A_t$, and take
$b\in D\setminus F\subseteq U$ with $F$ finite. For every earlier $s\in E$,
there is a point $a_s\in A_s\setminus U$. Such a point lies below $b$, hence in
$D$, and therefore in $F$. Distinct fibres supply distinct points, so there
are at most $|F|$ such $s$.

In any linear order whose strict initial segments are finite, the map
$t\mapsto|\{s\in E:s<t\}|$ is strictly increasing into $\N$. It is injective,
so $E$ is countable.
\end{proof}

\begin{proposition}\label[proposition]{prop:sep-fin}
For countable ordinals $\lambda$,
\[
 \sep(X_\lambda^\fin)=
 \begin{cases}
 n,&\lambda=n<\omega,\\
 \aleph_0,&\lambda=\omega,\\
 \cfr,&\omega<\lambda<\omega_1.
 \end{cases}
\]
\end{proposition}
\begin{proof}
For finite $n$, the space is discrete with $2^n$ points. The $n$ bit questions
separate it, and fewer binary questions produce fewer than $2^n$ answer words.
At $\lambda=\omega$, the countable family from \cref{prop:sep-down} remains
available, whereas no finite family separates an infinite set.

Now suppose $\omega<\lambda<\omega_1$. Partition $\B_\lambda$ into the ordered
fibres $A_u=C_u$, $u\in\two^\omega$. Each fibre contains at least two points.
For every $u$, a separating family must contain a set that splits $A_u$.
By \cref{lem:split-fibres}, one open set splits at most countably many fibres.
There are $\cfr$ fibres. Thus a separating family has cardinality at least
$\cfr$, since the union of $\kappa$ countable sets has size at most
$\max(\kappa,\aleph_0)$ for infinite $\kappa$. The family from
\cref{prop:sep-down} gives the matching upper bound.
\end{proof}

Consequently, among the countable values strictly above $\omega$, both
constructions keep cardinality and nonadaptive separation number fixed at
$\cfr$, while $\ps$ takes every such ordinal. The additional information is
not a cardinal-size effect; it records which queries can be made only after an
earlier transfinite prefix has become known.

\section{Verification, limitations, and further questions}

\subsection{What the accompanying computation checks}
The proofs above do not rely on an empirical extrapolation. The accompanying
Python program checks finite statements used to audit the reasoning:
\begin{enumerate}[label=(\alph*)]
\item the optimal comparison-query recurrence for chains of sizes up to $512$;
\item the untouched-child rule for every prefix and every initial-segment query
in all binary cubes of lengths $1$ through $10$;
\item the exact finite nonadaptive counts and the majority-side lower bound;
\item the finite pigeonhole assertion that more pairwise disjoint subcylinders
than deleted points leave at least one subcylinder untouched.
\end{enumerate}
The tests are exact integer and set computations, not floating-point experiments.
The saved output records the bounds and counts. They do not verify transfinite
induction, arbitrary open-set semantics, or the ZFC theorem in a proof assistant.
Those obligations are addressed by the written proofs, especially
\cref{lem:cylinders,lem:homogeneous,thm:fin-exact}.

Finite grids must be interpreted carefully. Replacing an infinite binary block
by a finite one makes the finite-deletion topology discrete on the whole finite
space. Such a model is useful for the local pigeonhole argument but cannot be
used as evidence that ordinal addition behaves like addition of finite query
counts. In particular, the limit-stage proof cannot be inferred from numerical
experiments.

\subsection{The crucial proof obligations}
The most important points for an independent check are the following.

\paragraph{Global legality of the queries.}
The set $D_s$ in \eqref{eq:seeker-cut} is globally downward closed. A local
coordinate condition on a known fibre is not silently assumed to be globally
open.

\paragraph{Nonempty limits.}
Every retained prefix has all possible extensions. Increasing prefixes unite
to a function on the union of their domains, and this function extends to the
whole ordinal. Thus the intersection of the corresponding cylinders remains a
full, nonempty cylinder.

\paragraph{Arbitrary refined opens.}
The finite-prefix lemma uses a basic neighborhood of just one point of an
arbitrary open set. It therefore handles arbitrary unions of basic sets, not
merely one downward set with one finite deletion.

\paragraph{Correct time accounting.}
In the refined topology, a single query may cost several prefix bits. Every
cost is finite within the active $\omega$-block. A strictly increasing sequence
of finite costs fills exactly that block at its limit, and cannot spend any
coordinate of the next block prematurely.

\paragraph{A genuine Hider strategy.}
The shortest-then-lexicographically-first homogeneous extension provides a
rule for each current position and query. No future Seeker strategy is consulted.
The last finite block uses a majority rule on tracked points. Thus the argument
proves Hider wins, not merely the failure of one proposed search algorithm.

\subsection{What is and is not settled}
Within the stated $T_0$ framework, \cref{thm:main} gives an affirmative answer
to \eqref{eq:problem}, with the stronger compact $T_1$ conclusion. It does not
settle a version requiring Hausdorff witnesses. In fact,
\cref{prop:compact} supplies an explicit obstruction to Hausdorffness for these
very spaces.

There is a concrete reason that the proof cannot simply be transferred to an
order topology. Two-sided open sets can cut independently inside many ordered
fibres. For example, the disjoint union of three nonempty open intervals in an
ordinary real line can split three disjoint, ordered intervals. This is
incompatible with the finite-predecessor conclusion of \cref{lem:split-fibres}.
Thus replacing the topology changes the query model, rather than merely
improving a separation axiom for free.

A natural next problem is to construct a Hausdorff space with value
$\omega^2+1$, or to find a different invariant that obstructs this. A successful
construction would need a replacement for the homogeneous-prefix lemma that
survives the additional open sets required for Hausdorff separation. The present
result identifies exactly what the one-sided ordered construction achieves,
and exactly which extra requirement is not addressed.

The source check is limited: the identified v3 explicitly lists the spectrum
question, and no later resolution was located in the searches performed for
this manuscript. This does not certify priority against all publications or
unpublished work. The proof is offered for mathematical checking on its own
terms, with all substantive steps supplied.

\section{Conclusion}

The ordinal-indexed binary cube, equipped with downward-closed queries, has
exact search length equal to its coordinate ordinal. The lower bound is not a
cardinality argument: it comes from retaining a full successor cylinder at each
turn and a full union cylinder at every limit. Passing to the cofinite refinement
preserves the exact value by finite-prefix homogeneity and $\omega$-block fusion,
while making the space $T_1$ and retaining compactness.

The resulting statement is
\[
 \boxed{\text{Every ordinal is the point-separating value of a compact }T_1
 \text{ space.}}
\]
In particular $X_{\omega^2+1}^\fin$ is the explicit witness requested by the
selected spectrum problem. The order-embedding characterization explains why
these constructions naturally belong to the study of order types and
transfinite notation: the game measures the least ordinal length of a
lexicographic binary code.

\appendix
\section{Strategy reference sheet}

\subsection*{Seeker, both topologies, game length $\lambda$}
At time $\eta$, retain the prefix determined by the previous answers, of length
$\eta$. Ask $D_s$ from \eqref{eq:seeker-cut}. Append $0$ after answer $1$, and
append $1$ after answer $0$. Unite the prefixes at limit stages. At time
$\lambda$, the cylinder is a singleton.

\subsection*{Hider, downward topology, any game length below $\lambda$}
Track a prefix $s$. If its left child is entirely contained in the query, keep
that child and answer $1$. Otherwise keep the right child and answer $0$.
Unite prefixes at limits. The prefix length equals game time, leaving a free
coordinate at every shorter stopping time.

\subsection*{Hider, finite-deletion topology, any game length below $\lambda$}
Write $\lambda=\omega\cdot\alpha+n$. Before time $\omega\cdot\alpha$, track a
prefix inside the active infinite block and choose a shortest nonempty finite
extension whose cylinder is homogeneous for the query. Such an extension exists
by \cref{lem:homogeneous}. Answer for its containing side. Take unions at all
limits. At time $\omega\cdot\alpha$, track all $2^n$ remaining points. For the
last finite block choose the larger side at each query, breaking ties in a fixed
way. Every stopping time before $\lambda$ leaves at least two tracked points.

\section{Reproducing the archive}
The archive contains this source, the compiled PDF, a Bib\TeX{} record of the
primary reference, the standard-library Python verification program, its saved
JSON output, and a separate proof-audit note. From the archive directory:
\begin{lstlisting}
python3 code/finite_checks.py --output verification/results.json
pdflatex -interaction=nonstopmode -halt-on-error ordinal_separation.tex
pdflatex -interaction=nonstopmode -halt-on-error ordinal_separation.tex
\end{lstlisting}
The PDF uses an embedded bibliography, so a Bib\TeX{} run is unnecessary.
A separate bibliography file, \texttt{references.bib}, is supplied for reuse. No third-party source PDF
or font file is redistributed.

\begin{thebibliography}{9}
\bibitem{CCS}
Lucas Chiozini, Tam\'as Csern\'ak, and Lajos Soukup,
\emph{Gamification of the $T_0$-pseudoweight via cut-and-choose games on
 topological spaces}, arXiv:2510.05754v3, submitted 29 May 2026,
PDF dated 1 June 2026.
The arXiv record uses the earlier title \emph{Cut-and-choose games in
topological spaces}.
\url{https://arxiv.org/abs/2510.05754v3};
versioned PDF: \url{https://arxiv.org/pdf/2510.05754v3}.
Definition~1.1, Problem~1.2, and the discussion of the ordinal spectrum.
Accessed 20 September 2026.
\end{thebibliography}
\end{document}
